\chapter{Conversor de divisas}
\label{chap:divisas}
Esta es la función de conversor de divisas de la aplicación.
Permite convertir un importe entre dólares, euros y yenes, y el resultado se recalcula mientras se escribe.

Se sigue la misma paleta de colores que en el resto de la aplicación.
La pantalla tiene una cabecera verde y cuatro fichas: divisa de origen, divisa de destino, importe y total.
El importe es un campo numérico que no permite escribir un signo negativo.

\begin{figure}[H]
  \centering
  \includegraphics[width=0.4\textwidth]{divisas}
\end{figure}

\section{Enumeración Divisa}
Cada divisa se define una sola vez con su nombre y su tasa respecto al euro.

\begin{lstlisting}
enum class Divisa(val nombre: String, val tasaPorEuro: Double) {
    DOLAR("Dólares", 1.08),
    EURO("Euros", 1.0),
    YEN("Yenes", 165.0)
}
\end{lstlisting}

Añadir una divisa nueva es añadir una línea.
Las tasas son fijas y de ejemplo, no se consultan a ningún servicio externo.

\section{Conversión}
Todas las conversiones pasan por el euro.
Se divide el importe entre la tasa de origen y se multiplica por la tasa de destino.

\begin{lstlisting}
fun convertir(importe: Double, origen: Divisa, destino: Divisa) =
    importe / origen.tasaPorEuro * destino.tasaPorEuro
\end{lstlisting}

Con una sola tasa por divisa se cubre cualquier pareja.

\section{Funcionamiento de la pantalla}
La divisa elegida en el origen se oculta en el destino, para evitar convertir una divisa en sí misma.
Cada vez que cambia el importe se llama a \texttt{calcularTotal()}.
Esta función cambia la coma por un punto y usa \texttt{toDoubleOrNull()}.
Si el texto no es un número válido, el total queda vacío y la aplicación no falla.
Si es válido, se muestra el total con dos decimales.